% Generated draft: requires mathematical and source review.
\documentclass[11pt]{article}
\usepackage[margin=25mm]{geometry}
\usepackage{fontspec}
\setmainfont{Latin Modern Roman}
\usepackage{amsmath,amssymb,mathtools,unicode-math}
\setmathfont{Latin Modern Math}
\usepackage{enumitem,needspace,xurl,hyperref}
\setlength{\parindent}{0pt}
\setlength{\parskip}{5pt}
\setlength{\emergencystretch}{3em}
\newcommand{\N}{\mathbb{N}}
\newcommand{\Z}{\mathbb{Z}}
\newcommand{\Q}{\mathbb{Q}}
\newcommand{\R}{\mathbb{R}}
\newcommand{\C}{\mathbb{C}}
\newcommand{\F}{\mathbb{F}}
\newcommand{\A}{\mathbb{A}}
\newcommand{\PP}{\mathbb{P}}
\newcommand{\E}{\mathbb{E}}
\newcommand{\eps}{\varepsilon}
\newcommand{\dd}{\,\mathrm{d}}
\begin{document}
\section*{UnsolvedMath ID: 30002764}
\textbf{Source problem code:} catalogue entry
\subsection{Definitions and mathematical statement}
Let \(G\) be a residually finite group, and let
\[
G=G_0\supseteq G_1\supseteq G_2\supseteq\cdots,
\qquad G_i\trianglelefteq G,
\qquad [G:G_i]<\infty,
\qquad \bigcap_{i\geq 0}G_i=\{e\}.
\]
Put \(Q_i=G/G_i\). For a finite matrix \(A\) over the rational group algebra \(\Q G\), let \(A_i\) be its image under the entrywise quotient map \(\Q G\to\Q Q_i\).

The supplied statement leaves the determinant convention and the meaning of normalization unstated. The following definitions make a possible formulation precise while retaining that uncertainty. For an \(m\times n\) matrix \(B\) over \(\Q H\), let \(T_B:\ell^2(H)^n\to\ell^2(H)^m\) be its left convolution operator, and put \(|T_B|=(T_B^*T_B)^{1/2}\). Write \(\mathcal N(H)\) for the group von Neumann algebra and use its canonical trace \(\tau_H(S)=\langle S\delta_e,\delta_e\rangle\). On matrices use the sum of the diagonal traces, without division by the matrix size. Define the spectral measure
\[
\mu_B(E)=\sum_{j=1}^n\left\langle 1_E(|T_B|)(\delta_e\otimes e_j),\delta_e\otimes e_j\right\rangle
\]
for Borel sets \(E\subseteq[0,\infty)\). Two conventions that must be distinguished are
\[
D_H^+(B)=\exp\left(\int_{(0,\infty)}\log t\,\mathrm d\mu_B(t)\right)
\]
and
\[
D_H^0(B)=
\begin{cases}
0,&\mu_B(\{0\})>0,\\
D_H^+(B),&\mu_B(\{0\})=0.
\end{cases}
\]
Here \(\exp(-\infty)=0\); the first convention omits the kernel, and in particular gives \(D_H^+(0)=1\).

For finite \(H\), the canonical trace corresponds to the ordinary matrix trace divided by \(|H|\). Thus, if \(s_1,\ldots,s_r\) are the positive singular values of the finite linear operator \(T_B\), counted with multiplicity, then
\[
D_H^+(B)=\left(\prod_{j=1}^r s_j\right)^{1/|H|},
\]
with empty product equal to \(1\). This specifies the finite-quotient normalization without applying a second normalization to a determinant already defined using the canonical trace.

\emph{Determinant approximation question, conditional on the determinant convention.} For every group, chain, and matrix as above, does
\[
D_G^{\star}(A)=\lim_{i\to\infty}D_{Q_i}^{\star}(A_i)
\]
hold, including existence of the limit, for the convention \(\star\) intended in the source? The supplied evidence does not explicitly determine whether \(\star=+\) or \(\star=0\), whether a kernel hypothesis is intended, or whether matrices must be square. Consequently these alternatives are not asserted to be equivalent or to be separately intended conjectures. The definitions accommodate rectangular matrices, but the intended range of matrix sizes requires source review. The original source fragment describes an analogue of L\"uck approximation for \(\ell^2\)-torsion; it supplies no chain-complex hypotheses from which a separate torsion statement can be specified.
\subsection{Short English statement}
The question asks whether the Fuglede--Kadison determinant of a rational group-ring matrix can be recovered from its images in every nested normal finite-quotient system with trivial intersection. The determinant convention, treatment of kernels, matrix dimensions, and finite-quotient normalization need confirmation. The supplied record labels the problem open; this draft does not independently verify that status.
\subsection{Sources}
\begin{itemize}
\item Local UnsolvedMath registry, supplied complete record for permanent ID \(30002764\), exact source code \texttt{catalogue entry}, titled ``Determinant Approximation for \(L^2\)-Torsion.'' The record belongs to the Oberwolfach Reports Open Problems collection and gives proposed year 2015; no proposer is named.
\item Supplied report citation: \emph{Geometric Topology} (2015), Oberwolfach Reports, pp.\ 187--233, DOI 10.4171/owr/2015/3. \url{https://doi.org/10.4171/owr/2015/3}. The supplied literature-source label is ``Grabowski, Oberwolfach Report''; this label is preserved as attribution and does not establish a named proposer.
\item Supplied literature reference: ``Lueck, approximation survey.'' \url{https://arxiv.org/abs/1501.07446}.
\item Literature assessment dated exactly \texttt{2026-08-21}: the supplied curation labels the problem open and describes determinant approximation for \(\ell^2\)-torsion as a major open problem, with known results for special residual systems but not the general claim. This is supplied dated curation, not an independently verified literature assessment; its time-sensitive status and cited claims require independent checking.
\item The supplied record assigns statement status \texttt{corrected\_verified} and reports comparison with the official Oberwolfach report. That is a supplied verification claim, not verification performed for this draft.
\end{itemize}
\end{document}
